% PHYS 331 — Exam 2 Formula Sheet
% Reorganized from E1 Formulae.pdf (two handwritten pages).
% Letter paper, two pages, intended for double-sided printing.
%
% EDITING GUIDE
%   1. Add formulas under the section that best fits them.
%   2. Insert new sections with \sectionheading{Section Name}.
%   3. The blank areas at the bottoms of the four columns are intentional.
%   4. After editing, check that the resulting PDF remains exactly two pages.
%   5. If needed, lower \fontsize{9.3}{11} on the line below or reduce
%      the gaps defined by \sectionheading, rather than changing paper size.

\documentclass[letterpaper,9pt]{extarticle}
\usepackage[letterpaper,top=0.4in,bottom=0.50in,left=0.53in,right=0.53in]{geometry}
\usepackage{amsmath,amssymb,mathtools}
\usepackage{braket}
\usepackage{microtype}
\usepackage{xcolor}
\usepackage{array}
\usepackage{ragged2e}
\pagestyle{empty}
\setlength{\parindent}{0pt}
\setlength{\parskip}{0pt}
\setlength{\abovedisplayskip}{4pt plus 1pt minus 1pt}
\setlength{\belowdisplayskip}{4pt plus 1pt minus 1pt}
\setlength{\abovedisplayshortskip}{2pt plus 1pt}
\setlength{\belowdisplayshortskip}{3pt plus 1pt}
\setlength{\arraycolsep}{3.2pt}
\renewcommand{\arraystretch}{1.05}
\definecolor{heading}{HTML}{26384A}
\definecolor{hairline}{HTML}{9BA7AE}
\newcommand{\pageheading}[2]{%
  {\large\bfseries\color{heading} PHYS 331 \enspace/\enspace Exam 2 Formula Sheet}%
  \hfill {\footnotesize\color{heading} #1 \quad\textbullet\quad #2}\par
  \vspace{3pt}\color{hairline}\hrule height 0.6pt\color{black}\vspace{9pt}}
\newcommand{\sectionheading}[1]{%
  \par\vspace{9pt}\noindent{\bfseries\color{heading} #1}\par
  \vspace{2pt}{\color{hairline}\hrule height 0.32pt}\vspace{5pt}}
\newcommand{\explain}[1]{\par{\footnotesize\color{heading}#1}\par\vspace{1pt}}
\newcommand{\columnbreakspace}{\vspace*{0pt}}

\begin{document}
\fontsize{9.3}{11}\selectfont
\pageheading{Avi Chetlin}{PHYS 331 E2}
\noindent
\begin{minipage}[t]{0.485\textwidth}
\sectionheading{Measurement and Expectation Values}
Probability of obtaining the eigenvalue associated with $\ket{a}$
\begin{equation*}
  p(a)=\bigl|\braket{a|\psi}\bigr|^2.
\end{equation*}
For an observable $\hat A$ with eigenvalues $a_n$,
\begin{equation*}
  \braket{\hat A}=\braket{\psi|\hat A|\psi}
  =\sum_n a_n p_n.
\end{equation*}

\sectionheading{Projectors and State Reduction}
Projection operator and corresponding outcome probability
\begin{align*}
  \hat P_a &= \ket{a}\!\bra{a}, \\
  p_a &= \braket{\psi|\hat P_a|\psi}
      =\braket{\psi|a}\braket{a|\psi}.
\end{align*}
Normalized state after projection onto the $n$th eigenspace
\begin{equation*}
  \ket{\psi'}=\frac{\hat P_n\ket{\psi}}
  {\sqrt{\braket{\psi|\hat P_n|\psi}}}.
\end{equation*}

\sectionheading{Operators in an Orthonormal Basis}
If $\{\ket{a_i}\}$ is an orthonormal basis, then
\begin{equation*}
  A_{ij}=\braket{a_i|\hat A|a_j},\qquad
  [\hat A]_{\{a_i\}}=(A_{ij}).
\end{equation*}
For example,
\begin{equation*}
  [\hat A]_{\{a_i\}}=
  \begin{pmatrix}
   \braket{a_1|\hat A|a_1} & \braket{a_1|\hat A|a_2} & \cdots\\
   \braket{a_2|\hat A|a_1} & \braket{a_2|\hat A|a_2} & \cdots\\
   \vdots & \vdots & \ddots
  \end{pmatrix}.
\end{equation*}
Conjugation identities
\begin{align*}
  \braket{\phi|\psi} &= \braket{\psi|\phi}^{*},\\
  \braket{\phi|\hat A|\psi}
  &= \braket{\psi|\hat A^{\dagger}|\phi}^{*}.
\end{align*}

\sectionheading{Quantum Information}
Bit flip/NOT gate
\begin{equation*}
\label{eq:1}
X = \sigma_z; X\ket{0} = \ket{1}, X\ket{1} = \ket{0}.
\end{equation*}
Phase gate
\begin{equation*}
\label{eq:2}
Z = \sigma_{z}; Z\ket{0} = \ket{0}, Z\ket{1} = -\ket{1}.
\end{equation*}
\begin{equation*}
Y = iZX = \sigma_y.
\end{equation*}
Hadamard gate
\begin{equation*}
H\ket{0} = \frac{1}{\sqrt{2}}(\ket{1} + \ket{0}), H\ket{1} = \frac{1}{\sqrt{2}}(-\ket{1}+\ket{0}).
\end{equation*}
\begin{equation*}
H \doteq \frac{1}{\sqrt{2}} \begin{pmatrix}
1  & 1 \\
1  & -1
 \end{pmatrix}.
\end{equation*}
Controlled gates



% ADD ADDITIONAL GENERAL-FORMALISM ENTRIES HERE.
\end{minipage}\hfill
\begin{minipage}[t]{0.485\textwidth}
\columnbreakspace
\sectionheading{Uncertainties}
Standard deviation of an observable
\begin{align*}
 \Delta A
 &=\sqrt{\braket{\hat A^2}-\braket{\hat A}^{2}}\\
 &=\sqrt{\braket{\psi|\hat A^2|\psi}
          -\braket{\psi|\hat A|\psi}^{2}}.
\end{align*}
Robertson uncertainty relation
\begin{equation*}
 \Delta A\,\Delta B\geq\frac12
 \left|\braket{[\hat A,\hat B]}\right|.
\end{equation*}

\sectionheading{Spin-$\tfrac12$ Matrices}
In the $S_z$ eigenbasis,
\begin{align*}
 \hat S_x&=\frac{\hbar}{2}\begin{pmatrix}0&1\\1&0\end{pmatrix},&
 \hat S_y&=\frac{\hbar}{2}\begin{pmatrix}0&-i\\i&0\end{pmatrix},\\[5pt]
 \hat S_z&=\frac{\hbar}{2}\begin{pmatrix}1&0\\0&-1\end{pmatrix}.&&
\end{align*}
Eigenvector matrices as written in the original notes
\begin{align*}
 P_{z\to x}&=\frac{1}{\sqrt{2}}
    \begin{pmatrix}1&1\\1&-1\end{pmatrix},\\[4pt]
 P_{z\to y}&=\frac{1}{\sqrt{2}}
    \begin{pmatrix}1&1\\i&-i\end{pmatrix},\\[4pt]
 P_{z\to z}&=\begin{pmatrix}1&0\\0&1\end{pmatrix}.
\end{align*}

\sectionheading{Spin Algebra}
\begin{equation*}
 \hat S^2=\hat S_x^2+\hat S_y^2+\hat S_z^2
 =\frac{3\hbar^2}{4}I.
\end{equation*}
For cyclic $(i,j,k)=(x,y,z)$,
\begin{equation*}
 [\hat S_i,\hat S_j]=i\hbar\hat S_k.
\end{equation*}

\sectionheading{Position and Momentum}
\begin{align*}
  \hat{x}\ket{x} &= x\ket{x} \\
  \ket{\psi} &= \sum_i \ket{x_i}\braket{x_i | \psi} \\
  \psi(x) &= \braket{x_i | \psi}.
\end{align*}
Probability
\begin{align*}
|\braket{x | \psi}|^2 &= |\psi(x)|^2 \\
\mathbb{P}(x_1 < x < x_{2})&=\int\limits_{x_1}^{x_2} |\psi(x)|^2 dx
\end{align*}
Momentum operator
\begin{align*}
\hat{p} = -i\hbar \frac{d}{dx}, \quad \langle \hat{p} \rangle = \int\limits_{-\infty}^{\infty} \psi^{*}(-i\hbar \frac{d\psi}{dx}) dx
\end{align*}
\begin{align*}
\left[ - \frac{\hbar^2}{2m}\frac{d^2}{dx^2} + V(x) \right]\varphi(x) = E\varphi(x)
\end{align*}

\end{minipage}

\newpage
\pageheading{Avi Chetlin}{PHYS 331 E2}
\noindent
\begin{minipage}[t]{0.485\textwidth}
\columnbreakspace
\sectionheading{Spin Along an Arbitrary Direction}
Unit vector in spherical coordinates
\begin{equation*}
 \hat{\mathbf n}=(\sin\theta\cos\phi,\,
 \sin\theta\sin\phi,\,\cos\theta).
\end{equation*}
Spin along $\hat{\mathbf n}$
\begin{align*}
 \hat S_n &=\mathbf{\hat S}\cdot\hat{\mathbf n}\\
 &=\hat S_x\sin\theta\cos\phi
 +\hat S_y\sin\theta\sin\phi
 +\hat S_z\cos\theta\\
 &=\frac{\hbar}{2}\begin{pmatrix}
 \cos\theta &e^{-i\phi}\sin\theta\\
 e^{i\phi}\sin\theta&-\cos\theta
 \end{pmatrix}.
\end{align*}
The eigenkets in the $S_z$ eigenbasis are
\begin{align*}
 \ket{+}_{\hat n}
 &=\cos\frac\theta2\ket{+}
 +e^{i\phi}\sin\frac\theta2\ket{-},\\[3pt]
 \ket{-}_{\hat n}
 &=\sin\frac\theta2\ket{+}
 -e^{i\phi}\cos\frac\theta2\ket{-}.
\end{align*}

\sectionheading{General Spin $s$}
For $m=-s,-s+1,\ldots,s$,
\begin{align*}
 \hat S^2\ket{s,m}&=s(s+1)\hbar^2\ket{s,m},\\
 \hat S_z\ket{s,m}&=m\hbar\ket{s,m}.
\end{align*}

\sectionheading{Square Wells}
Infinite
\begin{align*}
  \varphi_n(x) &= \sqrt{\frac{2}{L}} \sin(\frac{n\pi}{L}x), \quad E_n = \frac{n^2 \pi^2 \hbar^2}{2mL^2}, \\
  \psi(x) &= \sum\limits_n c_n\varphi_n(x) = \sum_n \left(\int_0^{L} \varphi^{*}_n(x)\psi(x) dx \right) \varphi_n(x)
\end{align*}
Finite
\begin{align*}
\varphi_{even} &= \begin{cases}
  Ae^{qx}, &x \leq -a \\
  D\cos kx, & -a < x < a \\
  Ae^{-qx}, & x \geq a,
\end{cases}, \quad k\tan(ka)=q \\
  z &\coloneq ka = \sqrt{\frac{2mEa^2}{\hbar^2}}, z_0 = \sqrt{\frac{2mV_{0}a^2}{\hbar^2}} \\
  z\tan z &= \sqrt{z_{0}^2 - z^2}
\end{align*}

\end{minipage}\hfill
\begin{minipage}[t]{0.485\textwidth}
\columnbreakspace
\sectionheading{Time Evolution}
Time-dependent Schr\"odinger equation
\begin{equation*}
 i\hbar\frac{\partial}{\partial t}\ket{\psi(t)}
 =\hat H\ket{\psi(t)}.
\end{equation*}
Energy eigenstates
\begin{equation*}
 \hat H\ket{E_n}=E_n\ket{E_n}.
\end{equation*}
Stationary-Hamiltonian expansion
\begin{align*}
 \ket{\psi}&=\sum_n c_n\ket{E_n},\\
 c_n(t)&=e^{-iE_n t/\hbar}c_n(0),\\
 \ket{\psi(t)}&=\sum_n c_n(0)e^{-iE_nt/\hbar}\ket{E_n}.
\end{align*}

\sectionheading{Spin in a Uniform Magnetic Field}
Magnetic moment and Hamiltonian
\begin{align*}
 \boldsymbol\mu&=g\frac{q}{2m}\mathbf{\hat S},&
 \boldsymbol\mu_e&\simeq-\frac{e}{m_e}\mathbf{\hat S},\\
 \hat H&=-\boldsymbol\mu\cdot\mathbf B.
\end{align*}
If $\mathbf B=B_0\hat{\mathbf z}$,
\begin{align*}
 \hat H&=\frac{eB_0}{m_e}\hat S_z
 =\omega_0\hat S_z,\\
 \omega_0&:=\frac{eB_0}{m_e},\\
 \hat H\ket{\pm}&=\pm\frac{\hbar\omega_0}{2}\ket{\pm}.
\end{align*}

\sectionheading{Complex and Trigonometric Identities}
\begin{align*}
 ie^{i\phi}&=e^{i(\phi+\pi/2)},&
 -e^{i\phi}&=e^{i(\phi+\pi)},\\
 e^{i\phi}+e^{-i\phi}&=2\cos\phi,\\
 e^{i\phi}-e^{-i\phi}&=2i\sin\phi,\\
 \cos\theta&=-\cos(\theta+\pi),\\
 \sin\theta&=-\sin(\theta+\pi).
\end{align*}

% ADD ADDITIONAL DYNAMICS AND COMPLEX-IDENTITY ENTRIES HERE.
\end{minipage}
\end{document}
